\documentclass[conference]{IEEEtran}

\usepackage[T1]{fontenc}
\usepackage[utf8]{inputenc}
\usepackage{amsmath,amssymb}
\usepackage{graphicx}
\usepackage{booktabs}
\usepackage{array}
\usepackage[table]{xcolor}
\usepackage{url}
\usepackage{microtype}

\graphicspath{{figures/}}
\setlength{\textfloatsep}{6pt plus 1pt minus 2pt}
\setlength{\floatsep}{5pt plus 1pt minus 2pt}
\setlength{\intextsep}{5pt plus 1pt minus 2pt}
\setlength{\abovecaptionskip}{3pt}
\setlength{\parskip}{0pt}

% =====================================================================
% FILL THESE IN DURING THE 20-MIN HACKATHON (single place to edit)
% =====================================================================
\newcommand{\TODO}{\textcolor{red}{\textbf{[TODO]}}}
\newcommand{\DSname}{RAW\_onsite}        % dataset folder name
\newcommand{\DSpmus}{\TODO}              % number of PMUs / CSV files
\newcommand{\DSbuses}{\TODO}             % observed bus numbers, e.g. 2,5,6,...
\newcommand{\DSrows}{\TODO}              % rows per CSV
\newcommand{\DSspan}{\TODO}              % minutes of data
\newcommand{\DSwindows}{\TODO}           % number of 30 s windows
\newcommand{\DSclasses}{\TODO}           % event classes present, e.g. 0,1,2,7
% --- metrics on the new dataset (leave [TODO] until scored) ---
\newcommand{\NewDetAcc}{\TODO}
\newcommand{\NewAbnRec}{\TODO}
\newcommand{\NewClsAcc}{\TODO}
\newcommand{\NewMacroF}{\TODO}
\newcommand{\NewLocTop}{\TODO}
\newcommand{\NewRuntime}{\TODO}          % wall-clock minutes
\newcommand{\NewFeatCount}{\TODO}        % feature count for this dataset
% --- features added on site (free text) ---
\newcommand{\FeaturesAdded}{\TODO\ (describe any feature added for this dataset)}
% =====================================================================

% figure placeholder that compiles without an image file; replace the
% whole \figslot{...} with \includegraphics[...]{figXX} when the plot is ready
\newcommand{\figslot}[2]{%
  \setlength{\fboxsep}{0pt}%
  \fbox{\begin{minipage}[c][#1][c]{0.96\columnwidth}\centering
    \color{gray}\footnotesize #2\\[2pt]\texttt{\textbackslash includegraphics here}
  \end{minipage}}}

\title{Physics-Informed Multi-PMU Event Detection and Localization:\\
On-Site Hackathon Report (SGSMA 2026, IEEE 39-Bus)}

\author{\IEEEauthorblockN{Jorge Luis Mayorga Taborda \and Yessica Africano Rodriguez}
\IEEEauthorblockA{Universidad de los Andes \\ SGSMA 2026 Synchrophasor Anomaly Detection Competition --- On-Site Round}}

\begin{document}
\maketitle

\begin{abstract}
We apply our pre-trained, CPU-only physics-informed ExtraTrees pipeline to the
on-site dataset \texttt{\DSname}. The method aligns the provided PMU streams,
builds baseline-relative electrical and data-quality features over fixed 30~s
windows, and routes each window through hierarchical detection, event
classification, and typed localization. This report summarizes the dataset, the
feature set (including features added on site), and the detection,
classification, and localization results.
\end{abstract}

\section{Method (unchanged from the submitted pipeline)}
The deployed model is bus-agnostic: PMUs are discovered from file names
(\texttt{Bus<N>\_*.csv}), so the same artifact runs on a new bus placement
without retraining. Each PMU stream is sorted by \texttt{TIMESTAMP} and split
into fixed 30~s windows. Per window we compute baseline-relative magnitude,
angle, frequency/ROCOF, derivative, robust $z$-score, sustained-anomaly,
data-quality, power-proxy, and graph-temporal features
($\approx\!45$k columns). A hierarchical ExtraTrees stack then predicts:
(i) abnormal vs.\ normal, (ii) the event class $\{0\text{--}8\}$ via a physical
classifier with physics gates for bad-data (7) and missing (5/6), and
(iii) the origin through typed BUS/LINE/PMU localizers. Output is the official
4-column CSV (\texttt{TIMESTAMP, Bus, Predicted\_Event, Predicted\_Location}).

\section{On-site dataset \texttt{\DSname}}
The provided package contains \DSbuses\ observed buses (\DSpmus\ PMU CSVs),
$\approx$\,\DSrows\ rows each spanning $\approx$\,\DSspan\ minutes at 30~fps,
yielding \DSwindows\ analysis windows. Event classes observed in the labels:
\DSclasses. Timestamps are aligned by row index across PMUs (an outer merge on
\texttt{TIMESTAMP} is avoided because of sub-ns drift between files).

\section{Features}
Table~\ref{tab:feat} lists the feature families used. For this dataset the
working feature count is \NewFeatCount.
\textbf{Added on site:} \FeaturesAdded.

\begin{table}[t]
\centering
\caption{Feature families fed to the ExtraTrees stack.}
\label{tab:feat}
\renewcommand{\arraystretch}{1.05}
\footnotesize
\begin{tabular}{@{}ll@{}}
\toprule
\textbf{Family} & \textbf{What it captures} \\
\midrule
Segment statistics & onset / transient / settling magnitude \\
Temporal deltas \& derivatives & abruptness, ramps, post-event shift \\
Robust $z$-score & scale-normalized abnormality \\
Sustained anomaly & duration above threshold (event vs.\ spike) \\
Data quality & NaNs, gaps, missing-run length \\
Single-signal / bad-data score & isolated channel faults \\
Power proxy ($P,Q$) & rough active/reactive movement \\
Graph-temporal & PMU-pair spread weighted by elec.\ distance \\
\bottomrule
\end{tabular}
\end{table}

\section{Results on \texttt{\DSname}}
Table~\ref{tab:res} reports detection, classification, and localization on the
on-site data, next to the submitted-model baselines for reference. Confusion
matrices and accuracy curves are shown in
Figs.~\ref{fig:cm}--\ref{fig:acc}.

\begin{table}[t]
\centering
\caption{Metrics: submitted baselines vs.\ on-site dataset.}
\label{tab:res}
\renewcommand{\arraystretch}{1.15}
\footnotesize
\begin{tabular}{@{}lccc@{}}
\toprule
\textbf{Metric} & \textbf{RAW0001} & \textbf{SIM} & \textbf{\DSname} \\
\midrule
Detection accuracy            & 1.0000 & 0.9693 & \NewDetAcc \\
Abnormal recall               & 1.0000 & 0.9860 & \NewAbnRec \\
Event classification accuracy & 1.0000 & 0.9673 & \NewClsAcc \\
Event classification macro-F1 & 1.0000 & 0.9598 & \NewMacroF \\
Localization Top-1 accuracy   & 0.8333 & 0.8524 & \NewLocTop \\
\bottomrule
\end{tabular}
\end{table}

\begin{figure}[t]
\centering
\figslot{6.2cm}{Detection / event confusion matrix on \DSname}
\caption{Confusion matrix on the on-site dataset. Replace the box with
\texttt{fig\_cm.pdf} once generated.}
\label{fig:cm}
\end{figure}

\begin{figure}[t]
\centering
\figslot{6.2cm}{Per-class precision / recall / F1 and accuracy bars}
\caption{Per-class metrics and overall accuracy on \DSname.}
\label{fig:acc}
\end{figure}

\section{Efficiency}
The runtime is CPU-only with \textbf{0 trainable neural parameters} and a model
bundle under \textbf{1~MB}. Inference on \texttt{\DSname} took \NewRuntime\
minutes on a laptop CPU (macOS). No GPU, no feature scaling; \texttt{NaN}s are
handled by a median imputer inside each model.

\section{Discussion}
The pipeline transfers without retraining because features are
baseline-relative and bus-agnostic. Expected sensitivities on unseen data:
(i) localization Top-1 depends on how many candidate buses are observable;
(ii) Event~8 (unknown) is used when no compact signature matches, which trades
recall on rare classes for precision; (iii) bad-data (7) and missing (5/6) are
decided by physics gates before the physical classifier, so communication
artifacts do not contaminate physical-event labels.

\end{document}
